
\documentclass[11pt,a4paper]{article}

\usepackage[T1]{fontenc}
\usepackage[utf8]{inputenc}
\usepackage[ngerman]{babel}
\usepackage[a4paper,margin=25mm]{geometry}
\usepackage{amsmath,amssymb}
\usepackage{helvet}
\usepackage{sfmath}
\renewcommand{\familydefault}{\sfdefault}
\usepackage{microtype}
\usepackage[hidelinks]{hyperref}
\usepackage{lastpage}
\usepackage{fancyhdr}
\usepackage{parskip}

\pagestyle{fancy}
\fancyhf{}
\fancyfoot[C]{\thepage\ (\pageref{LastPage})}
\renewcommand{\headrulewidth}{0pt}
\renewcommand{\footrulewidth}{0pt}

\hypersetup{
  pdftitle={MoE-Glossar: Interstitium},
  pdfauthor={Wolfgang Tscheu}
}

\begin{document}

\vspace*{-1.8cm}

\noindent
\begin{minipage}[t]{0.66\textwidth}
{\LARGE\bfseries Interstitium}
\end{minipage}
\hfill
\begin{minipage}[t]{0.30\textwidth}
\raggedleft
{\small MoE-Glossareintrag}\\
{\small Version 261029-1}
\end{minipage}

\vspace{1.4em}

\section*{Definition}

Ein Interstitium ist der räumliche Zwischenraum,
der durch die Oberflächen benachbarter, einander
berührender QUBLES begrenzt wird.

Die geometrische Gestalt des Interstitiums ergibt
sich aus der räumlichen Anordnung und den
Kontaktlagen der beteiligten QUBLES.

\section*{Tetraederähnliches Interstitium}

Vier einander paarweise berührende Kugeln
begrenzen zwischen ihren Oberflächen ein
konkaves, tetraederähnliches Interstitium.

Die Mittelpunkte der vier Kugeln bilden
die Eckpunkte eines Tetraeders.

Die Begrenzungsflächen des Interstitiums
sind Ausschnitte der gekrümmten
Kugeloberflächen.

In der klassischen Kugelpackungsgeometrie
wird dieser Zwischenraum als
\emph{Tetraederlücke} bezeichnet.

\section*{Oktaederähnliches Interstitium}

Sechs entsprechend angeordnete Kugeln
begrenzen zwischen ihren Oberflächen ein
konkaves, oktaederähnliches Interstitium.

Vier Kugelmittelpunkte bilden ein Quadrat.
Je ein weiterer Kugelmittelpunkt liegt
oberhalb und unterhalb der Quadratsebene.

Die sechs Mittelpunkte bilden die
Eckpunkte eines Oktaeders mit acht
dreieckigen Flächen.

Die Begrenzungsflächen des Interstitiums
werden durch die gekrümmten Oberflächen
der sechs Kugeln gebildet.

In der klassischen Kugelpackungsgeometrie
wird dieser Zwischenraum als
\emph{Oktaederlücke} bezeichnet.

\section*{Interstitien in QUBLE-Packungen}

Die räumliche Anordnung der QUBLES bestimmt
die Lage, Gestalt und Verbindung der
Interstitien innerhalb einer Packung.

Die Kontaktlagen und die Verformung der
QUBLES bestimmen die konkrete Ausbildung
der Zwischenräume.

Die jeweiligen Packungsordnungen werden
in den Glossareinträgen
\emph{SDQP / Superdichte QUBLE-Packung}
und
\emph{DQP / Dichte QUBLE-Packung}
beschrieben.

\vspace{0.9em}
\hrule
\vspace{0.45em}

{\footnotesize
MoE-Glossar:
\emph{QUBLE / Quantelung},
\emph{SDQP / Superdichte QUBLE-Packung},
\emph{DQP / Dichte QUBLE-Packung}.
}

\end{document}
